\chapter{メタアナリシス}\label{ch:15}

\begin{goalbox}
\begin{itemize}
\item 固定効果モデルと変量効果モデルを式で書き，推定対象の違いを説明できる．
\item 逆分散重み付き推定量を最尤法として導出し，その期待値・分散を求められる．
\item Cochran の $Q$ 統計量の期待値を導出し，DerSimonian--Laird のモーメント推定量 $\hat\tau^2$ と $I^2$ を求められる．
\item 変量効果モデルでの $\mu$ の推定量と信頼区間を計算し，フォレストプロットを読める．
\end{itemize}
\end{goalbox}

\begin{exambox}
\begin{itemize}
\item \kakomon{2021}{4}：脂質異常症治療薬の4試験の対数リスク比．変量効果モデル $Y_k=\mu+\delta_k+\varepsilon_k$ の周辺分布 $\Normal(\mu,\sigma_k^2+\tau^2)$，
      対数尤度，$\hat\mu=\sum\frac{y_k}{\sigma_k^2+\hat\tau^2}\big/\sum\frac{1}{\sigma_k^2+\hat\tau^2}$，$\hat\tau^2=0.013$ での数値（$\hat\mu\approx0.008$），$\tau^2=0.014$ を既知とした $\V[\hat\mu]$ と95\%CI．
\item \kakomon{2023}{3}：高脂血症治療薬の筋肉関連有害事象の4試験の対数オッズ比．固定効果の対数尤度と MLE，その期待値・分散，
      $Q$ の分解と $\E[Q]$，$\tau^2$ のモーメント推定量，$\sigma_k^2=\sigma^2$ の下での $\tau^2/(\tau^2+\sigma^2)$ の推定量 $(Q-K+1)/Q$，$Q=7.15$ で $0.58$．
\end{itemize}
\end{exambox}

\flowline{試験ごとの\\推定値と分散}{共通効果 $\theta$\\平均効果 $\mu$}{逆分散重み\\DL 法}{正規近似\\$\sigma_k^2$ 既知}{$Q$，$\hat\tau^2$，$I^2$}{異質性と\\一般化可能性}

\section{問題設定}\label{sec:15-setting}
$K$ 個の研究 $k=1,\dots,K$ から，効果の推定値 $Y_k$（対数 RR，対数 OR，対数 HR，平均差など）とその分散 $\sigma_k^2$ が得られる．
$\sigma_k^2$ は各研究の標本サイズから十分精度よく推定されているとして既知とみなし，重みを $w_k=1/\sigma_k^2$ とする．
比の指標は対数スケールで統合し，結果を指数変換して示す．

\section{固定効果モデル}\label{sec:15-fixed}
\Hissu
\begin{teigi}[固定効果（共通効果）モデル]\label{def:15-fe}
$Y_k\sim\Normal(\theta,\sigma_k^2)$，$k=1,\dots,K$ 独立．すべての研究が共通の真の効果 $\theta$ をもつ．
\end{teigi}
\begin{meidai}[\kakomon{2023}{3}〔1〕〔2〕]\label{prop:15-fe}
対数尤度 $l(\theta)=-\frac12\sum_k\{\log(2\pi\sigma_k^2)+w_k(y_k-\theta)^2\}$ を最大化して
\begin{equation}
\hat\theta=\frac{\sum_kw_kY_k}{\sum_kw_k},\qquad \E[\hat\theta]=\theta,\qquad \V[\hat\theta]=\frac{1}{\sum_kw_k}.\label{eq:15-fe}
\end{equation}
\end{meidai}
\begin{proof}
$\partial l/\partial\theta=\sum w_k(y_k-\theta)=0$ より $\hat\theta$．
$\V[\hat\theta]=\sum w_k^2\sigma_k^2/(\sum w_k)^2=\sum w_k/(\sum w_k)^2$．
\end{proof}
逆分散重み付き平均は，$\theta$ の線形不偏推定量の中で分散最小である．精度の高い（大きな）研究ほど重みが大きい．

\section{異質性と \texorpdfstring{$Q$}{Q} 統計量}\label{sec:15-q}
\Hissu
\begin{teigi}[Cochran の $Q$]\label{def:15-q}
$Q=\sum_kw_k(Y_k-\hat\theta)^2$，$\hat\theta$ は固定効果推定量．
\end{teigi}
固定効果モデルの下で $Q\sim\chi^2_{K-1}$ であり，均一性（$\theta_1=\dots=\theta_K$）の検定に用いる．
\begin{teiri}[$Q$ の期待値（\kakomon{2023}{3}〔3〕）]\label{thm:15-EQ}
任意の $\theta$ について $Q=\sum w_k(Y_k-\theta)^2-\sum w_k(\hat\theta-\theta)^2$ であり，$Y_k$ が独立で $\E[Y_k]=\theta$（変量効果モデルでは $\theta$ を $\mu$ と読む）のとき
\begin{equation}
\E[Q]=\sum_kw_k\V[Y_k]-\frac{\sum_kw_k^2\V[Y_k]}{\sum_kw_k}.\label{eq:15-EQ}
\end{equation}
\end{teiri}
\begin{proof}
$\sum w_k(Y_k-\theta)^2=\sum w_k\{(Y_k-\hat\theta)+(\hat\theta-\theta)\}^2=Q+2(\hat\theta-\theta)\sum w_k(Y_k-\hat\theta)+(\hat\theta-\theta)^2\sum w_k$ で，
中央の項は $\sum w_k(Y_k-\hat\theta)=0$ より消える．期待値をとると
$\E[\sum w_k(Y_k-\theta)^2]=\sum w_k\V[Y_k]$，$\E[(\hat\theta-\theta)^2]=\V[\hat\theta]=\sum w_k^2\V[Y_k]/(\sum w_k)^2$．
\end{proof}
固定効果モデルなら $\V[Y_k]=1/w_k$ で $\E[Q]=K-1$ である．

\section{変量効果モデル}\label{sec:15-random}
\Hinshutsu
\begin{teigi}[変量効果モデル]\label{def:15-re}
$Y_k=\theta_k+\varepsilon_k$，$\theta_k=\mu+\delta_k$，$\varepsilon_k\sim\Normal(0,\sigma_k^2)$，$\delta_k\sim\Normal(0,\tau^2)$，すべて独立．
$\mu$ は研究間の真の効果の分布の平均，$\tau^2$ は研究間分散．
\end{teigi}
\begin{meidai}[\kakomon{2021}{4}]\label{prop:15-re}
(1) $Y_k\sim\Normal(\mu,\sigma_k^2+\tau^2)$．\\
(2) 対数尤度 $\log L(\mu,\tau^2)=-\frac12\sum_k\bigl\{\log2\pi+\log(\sigma_k^2+\tau^2)+\frac{(y_k-\mu)^2}{\sigma_k^2+\tau^2}\bigr\}$．\\
(3) $\tau^2$ を与えたときの MLE と分散は，$w_k^*=1/(\sigma_k^2+\tau^2)$ として
\begin{equation}
\hat\mu=\frac{\sum_kw_k^*Y_k}{\sum_kw_k^*},\qquad \V[\hat\mu]=\frac{1}{\sum_kw_k^*}.\label{eq:15-re}
\end{equation}
\end{meidai}
\begin{proof}
(1) 正規分布の和の再生性，$\E[Y_k]=\mu$，$\V[Y_k]=\tau^2+\sigma_k^2$．(3) $\partial\log L/\partial\mu=\sum w_k^*(y_k-\mu)=0$．分散は\cref{prop:15-fe}と同じ計算．
\end{proof}
実際には $\tau^2$ を推定値 $\hat\tau^2$ で置き換える．$\tau^2$ の推定の不確かさは $\V[\hat\mu]$ に含まれないので，研究数が少ないと区間は狭すぎる傾向がある．

\subsection{DerSimonian--Laird 推定量}\label{subsec:15-dl}
\begin{teiri}[モーメント推定量（\kakomon{2023}{3}〔4〕，\cite{dl1986}）]\label{thm:15-dl}
変量効果モデルで $\V[Y_k]=\tau^2+1/w_k$ を\cref{eq:15-EQ}に代入すると
\[
\E[Q]=K-1+\tau^2\Bigl(\sum_kw_k-\frac{\sum_kw_k^2}{\sum_kw_k}\Bigr).
\]
$\E[Q]=Q$ とおいて
\begin{equation}
\hat\tau^2_{\mathrm{DL}}=\max\left\{0,\ \frac{Q-(K-1)}{\sum_kw_k-\sum_kw_k^2/\sum_kw_k}\right\}.\label{eq:15-dl}
\end{equation}
\end{teiri}
\begin{proof}
$\sum w_k(\tau^2+1/w_k)=\tau^2\sum w_k+K$，$\sum w_k^2(\tau^2+1/w_k)/\sum w_k=\tau^2\sum w_k^2/\sum w_k+1$ を\cref{eq:15-EQ}に代入する．
負になる場合は0に切り詰める．
\end{proof}
DL 推定量は反復不要で広く使われるが，ML や REML の推定値とは一般に異なる
（\kakomon{2021}{4} のデータでは $\hat\tau^2_{\mathrm{ML}}=0.013$ に対し $\hat\tau^2_{\mathrm{DL}}\approx0.032$）．

\subsection{\texorpdfstring{$I^2$}{I²}}\label{subsec:15-i2}
全研究で $\sigma_k^2=\sigma^2$（$w_k=w$）なら $\sum w-\sum w^2/\sum w=w(K-1)$ で $\hat\tau^2=\sigma^2\{Q-(K-1)\}/(K-1)$，これより
\begin{equation}
\widehat{\Bigl(\frac{\tau^2}{\tau^2+\sigma^2}\Bigr)}=\frac{Q-(K-1)}{Q}=:I^2\label{eq:15-i2}
\end{equation}
（\kakomon{2023}{3}〔5〕）．$I^2$ は「観測された推定値のばらつきのうち，偶然（研究内の誤差）ではなく研究間の真の差による部分の割合」と解釈される．ただしこれは一般には直感的な解釈であり，上のように研究内分散が等しい場合に限って $\tau^2/(\tau^2+\sigma^2)$ という分散比として導出できる（研究内分散が異なる場合は，$\sigma^2$ を「典型的な」研究内分散に置き換えた近似的な意味をもつ）．$I^2$ は
負なら0とする．目安として25\%，50\%，75\% がそれぞれ低・中・高の異質性とされる \cite{higgins2003}（$I^2$ の定義は \cite{higgins2002}）．
$Q=7.15$，$K=4$ なら $I^2=(7.15-3)/7.15=0.58$．

\section{推定対象の違い}\label{sec:15-target}
\begin{itemize}
\item 固定効果：「全研究に共通の効果 $\theta$」．異質性があるときは「各研究の効果の，精度による重み付き平均」を推定していると解釈できる．
\item 変量効果：「研究間で変動する真の効果の分布の平均 $\mu$」．研究を母集団からの標本とみなす．
\item 変量効果の重み $1/(\sigma_k^2+\tau^2)$ は，$\tau^2$ が大きいほど研究間で均等に近づき，小さな研究の相対的な重みが大きくなる．
\item 変量効果の CI は固定効果より広い（研究間変動を反映）．$\tau^2=0$ なら両者は一致する．
\end{itemize}
\begin{supbox}[補足：予測区間（出典：\cite{higgins2009}）]
新しい研究（あるいは新しい臨床現場）での真の効果の範囲を示すには，予測区間 $\hat\mu\pm t_{K-2}(0.025)\sqrt{\hat\tau^2+\V[\hat\mu]}$ を用いる．
$\mu$ の信頼区間が0を含まなくても，予測区間が0を含むことはよくある．
\end{supbox}

\section{フォレストプロット}\label{sec:15-forest}
各研究の点推定値と95\%CI を横棒で，重みを四角形の大きさで示し，最下段に統合推定値をひし形で示す図をフォレストプロットという．
横軸は通常，比の指標の対数目盛で，縦線が「効果なし」（RR $=1$）を表す．
\begin{figure}[htbp]
\centering
\begin{tikzpicture}[x=40mm,y=6mm]
\draw[gray] (0,0.3)--(0,-6.6);
\foreach \x/\l in {-1/-1.0,-0.5/-0.5,0/0,0.5/0.5}{\draw (\x,-6.6)--(\x,-6.75) node[below,font=\scriptsize]{\l};}
\node[font=\scriptsize] at (-0.25,-7.6) {対数オッズ比};
% studies: Y, se
\foreach \y/\m/\s/\w/\name in {-1/-0.40/0.2000/1.2/研究1,-2/-0.05/0.1414/1.7/研究2,-3/-0.60/0.3000/0.8/研究3,-4/0.15/0.1732/1.4/研究4}{
  \draw[thick] ({\m-1.96*\s},\y)--({\m+1.96*\s},\y);
  \fill[bkblue] ({\m-0.02*\w},{\y-0.133*\w})rectangle({\m+0.02*\w},{\y+0.133*\w});
  \node[anchor=east,font=\scriptsize] at (-1.25,\y) {\name};}
\fill[bkorange] (-0.298,-5)--(-0.1186,-4.7)--(0.0607,-5)--(-0.1186,-5.3)--cycle;
\node[anchor=east,font=\scriptsize] at (-1.25,-5) {固定効果};
\fill[bkred!70] (-0.4614,-6)--(-0.1684,-5.7)--(0.1247,-6)--(-0.1684,-6.3)--cycle;
\node[anchor=east,font=\scriptsize] at (-1.25,-6) {変量効果（DL）};
\end{tikzpicture}
\caption{\cref{reidai:15-1}のフォレストプロット}\label{fig:15-forest}
\end{figure}

\section{仮定}\label{sec:15-assump}
\begin{itemize}
\item 各研究の推定値の正規近似と，$\sigma_k^2$ を既知とみなせること（小さな研究やイベントの稀な研究では不正確）．
\item 研究の独立性．
\item 変量効果モデル：真の効果が正規分布に従い，研究の大きさと無関係であること．
\item 選択された研究が偏っていないこと（出版バイアス）．
\end{itemize}
\begin{supbox}[補足：出版バイアス]
有意な結果の研究ほど出版されやすいと，統合推定値は効果を過大評価する．
横軸に効果，縦軸に精度（$1/\SE$）をとったファネルプロットの非対称性で検討する．
\end{supbox}

\section{結果の解釈}\label{sec:15-interp}
\begin{itemize}
\item \kakomon{2021}{4}：$\tau^2=0.014$（既知）とすると $\hat\mu\approx0.007$，$\SE\approx0.105$，95\%CI $[-0.198,0.212]$（対数RR；$\hat\tau^2=0.013$ では $\hat\mu\approx0.008$）．高用量と標準用量で心血管以外の死亡リスクに差があるとは言えない．
\item $Q$ 検定は研究数が少ないと検出力が低い．有意でなくても異質性がないとは言えないので，$I^2$ と $\hat\tau^2$ で程度を示す．
\item 異質性が大きいときは統合値だけでなく，その原因（対象集団，用量，追跡期間）を部分集団解析やメタ回帰で探索する．
\end{itemize}

\section{手法選択}\label{sec:15-choice}
\Youchui $Q$ 検定の結果や $I^2$ の大きさだけで固定効果か変量効果かを機械的に選ばない．まず推定対象を決める：「全研究に共通の効果（または研究の精度で重み付けた平均）」を知りたいなら固定効果，「研究間で変動する真の効果の分布の平均」を知りたいなら変量効果である．そのうえで，以下を目安とする．
\begin{itemize}
\item 研究が同一のプロトコルで行われ異質性が小さい → 固定効果（MH 法，Peto 法を含む）．
\item 異なる集団・条件の研究を統合し平均的効果を知りたい → 変量効果（DL，REML）．
\item 2値データでイベントが稀 → Peto 法や MH 法（逆分散法はゼロセルに弱い）．
\end{itemize}

\section{よくある誤り}\label{sec:15-err}
\begin{warnbox}
\begin{itemize}
\item 比の指標を元のスケールで平均する．
\item 変量効果の重みに $1/\sigma_k^2$ を使う（正しくは $1/(\sigma_k^2+\tau^2)$）．
\item $Q$ の自由度を $K$ とする（正しくは $K-1$）．$\hat\tau^2<0$ をそのまま使う．
\item $Q$ 検定が非有意なら固定効果でよいと機械的に判断する．
\end{itemize}
\end{warnbox}

\section{数理統計との接続}\label{sec:15-link}
\begin{linkbox}
\begin{itemize}
\item 逆分散重み付き平均が線形不偏推定量の中で最良であることは『現代数理統計学の基礎』第6章演習（問9）で，固定効果推定量の原型である．
\item 変量効果モデルは正規--正規の階層モデル（同書 4.2.4項，6.4.1項の経験ベイズ）であり，各研究の $\theta_k$ の事後平均は $Y_k$ を $\hat\mu$ の方へ縮小した値になる．
\item $\E[Q]$ の計算は平方和の分解 $\sum(X_i-\mu)^2=\sum(X_i-\bar X)^2+n(\bar X-\mu)^2$ の重み付き版で，不偏分散の除数 $n-1$ の導出と同じ構造である．
\end{itemize}
\end{linkbox}

\section{例題}\label{sec:15-ex}
\begin{reidai}[固定効果と変量効果]\label{reidai:15-1}
4研究の対数オッズ比 $Y_k$ と分散 $\sigma_k^2$ が $(-0.40,0.04)$，$(-0.05,0.02)$，$(-0.60,0.09)$，$(0.15,0.03)$ である．
\begin{enumerate}[label=(\arabic*)]
\item 固定効果推定値とその95\%CI を求めよ．
\item $Q$ を求め，均一性を検定せよ（$\chi^2_3(0.05)=7.815$）．$I^2$ を求めよ．
\item DL 法で $\hat\tau^2$ を求めよ．
\item 変量効果推定値とその95\%CI を求め，(1)と比べよ．
\end{enumerate}
\end{reidai}

\begin{reidai}[{$\E[Q]$ の応用}]\label{reidai:15-2}
$K=5$ 研究がすべて同じ分散 $\sigma^2=0.05$ をもつ．$Q=12.0$ のとき，$\hat\tau^2_{\mathrm{DL}}$ と $I^2$ を求めよ．
\end{reidai}

\section{解答}\label{sec:15-sol}
\begin{kaitou}{reidai:15-1}
(1) $w=(25,50,11.11,33.33)$，$\sum w=119.44$，$\sum wY=-10-2.5-6.667+5=-14.17$．
$\hat\theta=-14.17/119.44=-0.119$，$\SE=1/\sqrt{119.44}=0.0915$，95\%CI $(-0.298,\,0.061)$（OR で $(0.74,1.06)$）．\\
(2) $Q=\sum w_k(Y_k+0.119)^2=1.980+0.235+2.575+2.405=7.19<7.815$．均一性は棄却されない．$I^2=(7.19-3)/7.19=0.58$．\\
(3) $\sum w^2=625+2500+123.5+1111.1=4359.6$，$\sum w-\sum w^2/\sum w=119.44-36.50=82.95$．$\hat\tau^2=(7.19-3)/82.95=0.0506$．\\
(4) $w^*=1/(\sigma_k^2+0.0506)=(11.04,14.17,7.11,12.41)$，$\sum w^*=44.74$，$\sum w^*Y=-7.53$．
$\hat\mu=-0.168$，$\SE=1/\sqrt{44.74}=0.150$，95\%CI $(-0.461,\,0.125)$．
変量効果では小さい研究（研究3）の相対的な重みが増えて点推定が負の側に移り，区間は研究間変動を反映して広くなった．
$Q$ 検定は有意でないが $I^2=58\%$ で中程度の異質性があり，研究数が少ないため検定の検出力が低いことに注意する．
\end{kaitou}

\begin{kaitou}{reidai:15-2}
$\hat\tau^2=\sigma^2\{Q-(K-1)\}/(K-1)=0.05\times8/4=0.10$．$I^2=(12-4)/12=0.667$（$=0.10/(0.10+0.05)$）．
\end{kaitou}

\begin{summarybox}
\begin{itemize}
\item 固定効果：$\hat\theta=\sum w_kY_k/\sum w_k$，$\V=1/\sum w_k$，$w_k=1/\sigma_k^2$．
\item $Q=\sum w_k(Y_k-\hat\theta)^2\sim\chi^2_{K-1}$（均一性の下）．$\E[Q]=\sum w\V[Y]-\sum w^2\V[Y]/\sum w$．
\item 変量効果：$Y_k\sim\Normal(\mu,\sigma_k^2+\tau^2)$，$\hat\mu=\sum w_k^*Y_k/\sum w_k^*$，$w_k^*=1/(\sigma_k^2+\tau^2)$，$\V=1/\sum w_k^*$．
\item DL：$\hat\tau^2=\max\{0,(Q-K+1)/(\sum w-\sum w^2/\sum w)\}$．$I^2=(Q-K+1)/Q$．
\item 固定効果は共通効果，変量効果は効果の分布の平均．変量効果の CI は広く，小研究の重みが相対的に大きい．
\end{itemize}
\end{summarybox}
